\documentclass[11pt]{article}

\usepackage[T1]{fontenc}
\usepackage[utf8]{inputenc}
\usepackage{times}
\usepackage{microtype}
\usepackage[a4paper,margin=1in]{geometry}
\usepackage{amsmath,amssymb}
\usepackage{booktabs}
\usepackage{tabularx}
\usepackage{array}
\usepackage{enumitem}
\usepackage{xcolor}
\usepackage[numbers,sort&compress]{natbib}
\usepackage{url}
\usepackage{hyperref}

\hypersetup{
  colorlinks=true,
  linkcolor=black,
  citecolor=blue,
  urlcolor=blue
}
\setlist{nosep,leftmargin=1.5em}
\renewcommand{\arraystretch}{1.18}
\newcommand{\method}{\textsc{MemImpact}}
\newcommand{\R}{\mathcal{R}}
\newcommand{\A}{\mathcal{A}}

\title{\textbf{\method: Planner-Level Counterfactual Memory\\Auditing for Selective Verification}\\[0.5em]
\large COMP7607A Natural Language Processing --- Group Project Proposal}
\author{
\begin{tabular}{ll}
\multicolumn{2}{c}{\textbf{Team Members}}\\
Cao Junjia & Student ID: 3036790730\\
Zhang Yankai & Student ID: 3036806682\\
Wu Hongyao & Student ID: 3036793196\\
Kwan Tsz Hong & Student ID: 3036792350\\
Wang Zicheng & Student ID: 3036790326
\end{tabular}}
\date{}

\begin{document}
\maketitle

\begin{abstract}
Persistent memory helps tool-using language-model agents reuse earlier information, but stale or misattributed memories can produce wrong actions. We study \method, which removes each retrieved memory in turn, re-runs the planner, and uses changes in the proposed action to rank memories for verification. An action flip measures behavioral dependence, not whether verification will improve the final outcome. Using Qwen, LangGraph, and ToolSandbox, we compare \method\ and a flip-only variant with Always Proceed, Always Verify, Static, and a literal argument-overlap baseline. The evaluation contains 120 held-out synthetic episodes with direct and indirect memory dependence, and two independently run verification-budget settings. Separate forced-verification trials identify memories whose checking improves execution. We evaluate task success, harmful execution, abstention, verification calls, and planner tokens. The existing results show useful corrections in some indirect cases but do not establish an overall reliability--cost advantage over simpler policies. The project therefore tests both the value and the limits of planner re-execution as a verification signal.
\end{abstract}

\section{Introduction}

Long-lived LLM agents increasingly rely on persistent memory to preserve user preferences, environment facts, and execution traces across sessions \cite{packer2023memgpt,xu2025amem}. Memory can improve personalization and long-horizon performance, but its usefulness is not guaranteed by semantic relevance or successful retrieval. Retrieved experiences can propagate errors, induce misaligned replay, or remain confidently asserted after their original context has changed \cite{xiong2026memory,kwon2026confidence}. Recent benchmarks similarly show a gap between recalling updated evidence and acting correctly on it \cite{chao2026stale,he2026memoryarena}.

A natural response is to verify memories before using them. However, verifying every retrieved item introduces model calls, tool use, latency, and potentially unnecessary user interruptions. The verification decision therefore depends not only on whether a memory appears reliable, but also on whether it materially affects the proposed action and whether the action's consequences justify additional checking.

Several lines of prior work address adjacent parts of this problem. Causal Memory Intervention (CMI) selects memories by estimated causal usefulness \cite{srivastava2026cmi}. MemAudit attributes already-observed harmful behavior to poisoned memories through counterfactual replay \cite{tan2026memaudit}. ConMem uses prospective counterfactual scenarios for write-time retention decisions \cite{zhang2026conmem}. Decision-aware memory cards rank context using counterfactual-inspired utility \cite{guan2026cicl}. Most directly, AMV-Guard performs action-conditioned verification by extracting premises of a fixed candidate action and applying calibrated value of information \cite{wang2026amv}. These studies establish that counterfactual analysis and selective verification are not new by themselves.

We instead focus on a narrower unresolved question: \textbf{can re-running the planner without a memory provide a useful online verification signal that is not available from static memory features or a fixed candidate action alone?} A candidate-conditioned audit asks which premises support the action already proposed. A planner-level audit asks whether the agent would have proposed a different action in the first place. The two signals need not coincide, especially when a memory changes destinations, constraints, tool arguments, or the plan itself.

\section{Project Goal and Research Questions}

\subsection{Goal}

The project evaluates whether planner-level action changes help select memories for verification before a tool action. It measures the benefit of the resulting checks and the full cost of producing the counterfactual signal. The goal is a conditional account of when planner re-execution is worthwhile, including cases in which simpler policies perform better.

\subsection{Research Questions}

\begin{enumerate}[label=\textbf{RQ\arabic*:},leftmargin=3.2em]
  \item \textbf{Impact versus verification benefit.} How often does a noise-adjusted action flip select a memory whose forced verification improves the executed outcome, and how often does it trigger an unnecessary check?
  \item \textbf{Reliability--cost trade-off.} How do task success, harmful execution, abstention, verification calls, and total planner tokens compare with budget-matched Static, literal argument-overlap, and flip-only policies, as well as unconstrained Always Verify?
  \item \textbf{Scope and failure modes.} How do outcomes differ between direct and indirect memory dependence and across memory conditions? How do planner variability and independent runs limit conclusions about the method?
\end{enumerate}

\subsection{Objectives}

\begin{enumerate}
  \item Implement a tool agent that receives fixed prior-session memory snapshots and records structured actions in a reversible sandbox.
  \item Define an empirical planner-level memory-impact measure based on paired rollouts and an estimate of ordinary planner variability.
  \item Rank positive-impact memories within a fixed verification budget, verify selected records, and re-plan when authoritative evidence refutes a memory.
  \item Construct a held-out evaluation suite with executable ground truth for actions and verification outcomes, including both direct-argument and indirect action-selection dependencies.
  \item Compare six policies on the same tasks and initial candidate plans while counting counterfactual audit calls as part of total computational cost.
\end{enumerate}

\section{Related Work}

\subsection{Persistent Memory and Memory Evaluation}

MemGPT introduced operating-system-inspired context and memory management \cite{packer2023memgpt}, while A-MEM dynamically organizes structured memories and their links \cite{xu2025amem}. LongMemEval evaluates information extraction, multi-session reasoning, temporal reasoning, knowledge updates, and abstention in conversational assistants \cite{wu2025longmemeval}. MemoryArena moves beyond recall-centric evaluation by coupling memory acquisition and subsequent action in interdependent multi-session tasks \cite{he2026memoryarena}. STALE specifically demonstrates that models often fail to revise invalidated beliefs and adapt downstream behavior \cite{chao2026stale}. These benchmarks motivate evaluation at the level of decisions, not retrieval alone.

Work on memory management further shows that agents tend to follow retrieved experiences, allowing both errors and misleading executions to propagate \cite{xiong2026memory}. Provenance-aware TierMem addresses a complementary cost problem by escalating from compressed summaries to raw logs when summary evidence is insufficient \cite{zhu2026tiermem}. Manufactured Confidence shows that consolidation can turn hedged statements into confident assertions and that a single load-bearing memory can drive consequential errors \cite{kwon2026confidence}. Our work retains provenance and uncertainty as inputs but focuses on measuring the planner's behavioral dependence on each retrieved memory.

\subsection{Counterfactual Memory Selection and Auditing}

CMI evaluates controlled memory interventions and selects memories based on causal usefulness in Causal-LoCoMo \cite{srivastava2026cmi}. MemAudit combines a counterfactual memory influence score with structural anomaly detection to identify poisoned records after harmful behavior has been observed \cite{tan2026memaudit}. ConMem evaluates candidate memories through synthesized future scenarios to make retention decisions \cite{zhang2026conmem}. CICL ranks files, traces, rules, and memories with counterfactual-inspired decision utility and compacts them into typed memory cards \cite{guan2026cicl}. These methods differ in their intervention target and stage of the memory lifecycle; none directly evaluates our proposed online gate based on planner re-execution for pre-action verification.

AMV-Guard is the closest prior work. It extracts premises supporting a candidate action, calibrates their error probability, and selects verification actions by bounded value of information \cite{wang2026amv}. Our study compares planner-level rollouts with a simpler heuristic based on literal action-argument overlap. This baseline tests one form of fixed-candidate inspection; it does not reproduce AMV-Guard's premise extraction or value-of-information estimator.

\begin{table}[t]
\centering
\small
\caption{Positioning of the proposed project relative to directly related work.}
\label{tab:related}
\begin{tabularx}{\textwidth}{>{\raggedright\arraybackslash}p{0.19\textwidth} >{\raggedright\arraybackslash}p{0.27\textwidth} >{\raggedright\arraybackslash}p{0.23\textwidth} >{\raggedright\arraybackslash}X}
\toprule
\textbf{Work} & \textbf{Primary objective} & \textbf{Intervention or decision object} & \textbf{Distinction of \method} \\
\midrule
CMI \cite{srivastava2026cmi} & Select useful memories for downstream answers & Memory inclusion and answer quality & Studies pre-action verification rather than memory selection alone \\
MemAudit \cite{tan2026memaudit} & Identify poisoned memories after harmful behavior & Post-hoc counterfactual attribution & Targets online prevention and ordinary stale/conflicting memories as well as adversarial ones \\
ConMem \cite{zhang2026conmem} & Improve write-time memory retention & Prospective synthetic scenarios & Evaluates an already retrieved memory at action time \\
CICL \cite{guan2026cicl} & Rank and compress decision-relevant context & Counterfactual-inspired context utility & Uses actual planner rollouts rather than only judge-estimated utility \\
AMV-Guard \cite{wang2026amv} & Verify premises of a fixed candidate action using VOI & Candidate-conditioned premise verification & Measures whether removing a memory changes the planner's proposed action \\
\midrule
\method\ (this project) & Allocate online verification by planner-level behavioral impact & Paired re-planning with and without each memory & Directly compares planner re-execution with static and candidate-conditioned signals \\
\bottomrule
\end{tabularx}
\end{table}

\section{Proposed Method}

\subsection{Problem Setup}

Let $x$ denote the current request and $\R=\{m_1,\ldots,m_k\}$ the fixed retrieved memory set. Each memory has a claim, entity, scope, timestamp, and provenance. The planner $\Pi$ proposes a structured tool action $a^{+}=\Pi(x,\R;\xi)$, where $\xi$ denotes run-to-run variation. We canonicalize tool names and arguments for comparison \cite{yao2023react}. The online policies do not receive the simulator's hidden ground truth or offline critical-memory labels.

The LangGraph workflow plans, audits memory removals, selects verification targets, verifies them against authoritative scenario records, replans when evidence refutes a claim, and executes at most one state-changing tool action in ToolSandbox. The verifier returns \texttt{confirmed}, \texttt{refuted}, or \texttt{unknown}. The planner receives corrective evidence on replanning; the implementation does not run a fresh memory-removal audit after that replan.

\subsection{Planner-Level Counterfactual Audit}

For each retrieved memory $m_i$, we construct a paired ablated rollout:
\begin{equation}
  a_i^{-}=\Pi\!\left(x,\R\setminus\{m_i\};\xi\right).
\end{equation}
The same seed or decoding configuration is used for the full and ablated runs where supported. To account for stochastic planning, we estimate the empirical action-flip rate over $S$ paired runs:
\begin{equation}
  \widehat{F}_i=\frac{1}{S}\sum_{s=1}^{S}
  \mathbf{1}\!\left[\A\!\left(a^{+,(s)}\right)\neq
  \A\!\left(a_i^{-,(s)}\right)\right],
\end{equation}
where $\A(\cdot)$ returns the canonical next action, and equivalence is checked using tool names and normalized arguments rather than raw text alone. When stochastic decoding requires $S>1$, pairwise comparisons among the existing full-memory runs provide a baseline flip rate $\widehat{F}_0$ without additional planner calls. We use the noise-adjusted signal
\begin{equation}
  \widetilde{F}_i=\max\left(0,\widehat{F}_i-\widehat{F}_0\right).
\end{equation}
The frozen runs use $S=3$ with the same decoding settings. Full-memory and ablated planning traces are shared across \method\ and flip-only within each run. These planning-only branches do not execute state-changing tools.

This is an empirical planner-level intervention, not a correctness label or a formal causal estimate over all unobserved factors. Ordinary planner variability can create or suppress apparent flips. Leave-one-out removal can also miss redundant or interacting memories; we discuss this limitation without claiming a joint-removal experiment in the core study.

\subsection{Selective Verification Policy}

The implemented priority score is
\begin{equation}
  V_i=\frac{L\,q_i\,\widetilde{F}_i}{c_i},
  \qquad
  q_i=\left(1+\min\left(\frac{\operatorname{age}_i}{365\text{ days}},1\right)\right)r_i,
\end{equation}
where $r_i=1$ when provenance is present and $1.25$ otherwise. $L$ is the episode's predefined consequence weight and $c_i=1$ is the unit cost of a memory-verification check. Neither is estimated by a separate model call in the core system. The gate selects positive-score memories in descending order, up to a cap of one or three checks; it has no additional learned or tuned threshold. Because $L$ and $c_i$ are constant across memories in an episode, they do not change which memories are selected under this rule. The comparison with flip-only chiefly tests how age and provenance weighting resolve score ties.

Selected memories are checked through the same authoritative verifier for every policy. Refuting evidence triggers a planner call with corrective evidence before execution. An unknown result for a task-critical memory can cause abstention. A replanned state-changing action that fails to use a verified value is blocked. Read-only contact or reminder lookup may also trigger replanning; execution is capped at one state-changing action and two replans. Verification outcomes are logged, but the core run does not perform another full MemImpact audit after verification.

\section{Experimental Plan}

\subsection{Environment and Memory Scenarios}

The primary environment is ToolSandbox \cite{lu2025toolsandbox}, augmented with synthetic prior-session memory snapshots and authoritative verification records. ToolSandbox provides stateful tool execution and dynamic trajectory evaluation. The study does not implement an open-ended memory-writing system or claim transfer to AppWorld \cite{trivedi2024appworld}.

The dataset contains single-decision episodes with synthetic prior-session memories and executable ground truth. Forty development episodes cover contact updates and reminder creation; 120 held-out episodes cover messaging and reminder deletion. Related variants of a base scenario stay in the same split. Core memory conditions include:
\begin{itemize}
  \item \textbf{Valid:} current and correctly scoped memories, including old but still valid records;
  \item \textbf{Stale:} previously valid memories superseded by newer authoritative records;
  \item \textbf{Conflicting:} memories inconsistent with another retrieved source;
  \item \textbf{Misattributed or over-generalized:} facts assigned to the wrong person or scope;
  \item \textbf{Irrelevant but similar:} semantically related memories that should not affect the decision.
\end{itemize}
Adversarial memories are outside the evaluated core matrix.

Each episode also varies the dependence type: a \emph{direct} memory supplies a tool argument, such as a phone number, whereas an \emph{indirect} memory changes which action is selected without appearing verbatim in its arguments, such as a rule specifying which contact channel to use. This factor separates argument overlap from action-selection influence. The held-out matrix is $2$ task families $\times$ $5$ memory conditions $\times$ $2$ consequence levels $\times$ $2$ dependence types $\times$ $3$ scenario variants, yielding 120 episodes. Consequence level varies independently of memory correctness.

\subsection{Baselines}

All methods share the same base agent, initial candidate action, retrieved memory set, tools, and candidate tasks.
\begin{enumerate}
  \item \textbf{Always Proceed:} trust all retrieved memories without verification;
  \item \textbf{Always Verify:} verify every retrieved candidate memory within the same verifier interface; this is a full-check reference, not a budget-matched selective policy;
  \item \textbf{Static Policy:} select the oldest memories in task-relevant scopes without counterfactual rollouts;
  \item \textbf{Literal argument-overlap policy:} select memories whose claims contain a value appearing in the initial action arguments, without re-running the planner.
\end{enumerate}
The literal overlap policy is an AMV-inspired, candidate-conditioned heuristic. It does not perform learned premise extraction or reproduce AMV-Guard's calibration and value-of-information procedure. Static and literal overlap use the same initial candidate, verifier, executor, and selective-verification budget as \method; the full-check policy verifies all retrieved memories independently of that cap.

\subsection{Ablation Studies}

The implemented ablation is flip-only, which ranks by $\widetilde{F}_i$ with the same three paired samples and verification budget as the full policy. It tests whether age and provenance weighting improve selection beyond action flips alone. We also analyze how full-memory planner variability affects the adjusted score. Joint-memory removal, adaptive sampling, and pre-audit triggers remain future work and are not reported as completed experiments.

\subsection{Evaluation Metrics}

\begin{itemize}
  \item \textbf{Primary outcomes:} end-to-end task success, harmful execution over all episodes, abstention rate, verification calls per episode, and total planner tokens per episode;
  \item \textbf{Identification:} precision and recall of selected memories against labels from separate offline forced-verification-and-replanning trials;
  \item \textbf{Secondary outcomes:} safe-proceed recall, proceed coverage, planner calls and attempts, planner latency, and direct/indirect subgroup results;
  \item \textbf{Counterfactual quality:} ordinary planner variability, noise-adjusted action-flip rates, and selected failure cases.
\end{itemize}
An offline label is positive only when forcing verification of one memory changes the shared initial failed or harmful outcome into a successful, non-harmful one. These labels are unavailable to the online policy. Action changes alone are not correctness labels, and offline label-generation tokens are reported separately. Online planner-token totals include initial planning, audit rollouts, retries, and replanning. Verification calls are counted separately. Planner latency does not include verification-tool execution, and we do not estimate monetary cost from incomplete pricing data.

\subsection{Controls and Statistical Analysis}

The study uses one frozen hosted Qwen model configuration and identical scenario fixtures within each run. The initial full-memory plan is shared across all six policies, and full-memory and ablated traces are shared by \method\ and flip-only within a run. Selective policies are evaluated with caps of one and three memory checks. The two budget configurations were run with independent model calls, so differences between them are descriptive and cannot be attributed solely to budget. Always Verify is an unconstrained full-check reference.

Policy differences within each run are summarized with paired bootstrap intervals clustered by base scenario. Subgroups with few independent base scenarios are descriptive. Offline critical-memory labels are generated after policy execution from separate forced-verification trials. The study does not use held-out outcomes to tune a decision threshold. Claims about additional strategies suggested by these results would require a new, untouched evaluation set.

\subsection{Reproducibility, Ethics, and Safety}

All users, memories, contacts, accounts, and transactions are synthetic. Experiments run in isolated sandboxes without real messages or external side effects. The repository contains configurations, prompts, scenario rules, traces, evaluation scripts, and a claim-to-artifact table linking reported results to logs. No human-subject data is used. The agent is not presented as providing a general safety guarantee.

\section{Expected Outcomes and Limitations}

The project delivers a working \method\ pipeline, a reproducible synthetic memory-verification evaluation, and an analysis of action flips, verification benefit, and computational cost. With a one-check cap, \method\ reaches 94.17\% success versus 84.17\% for flip-only; with a three-check cap, the corresponding rates are 93.33\% and 92.50\% in an independent run. Always Verify reaches 99.17\% and 100\% success in the two runs while using far fewer planner tokens. These observations support localized value in some indirect cases, but not a reliable overall reliability--cost advantage for planner re-execution.

Planner rollouts are expensive and vary even with fixed decoding settings. Leave-one-out removal can miss redundant or interacting memories, and subtraction of the full-memory flip baseline can suppress a useful signal. The score's age and provenance factors are heuristics rather than calibrated error probabilities; the consequence weight comes from scenario metadata and verification cost is fixed at one unit. The 120 held-out episodes contain only two task families and twelve base-scenario clusters, so subgroup estimates and external validity are limited. Authoritative verification and simulator-derived harm cover only encoded outcomes. The literal-overlap baseline is weaker than a full AMV-Guard implementation. We do not present the empirical action-flip score as a general safety guarantee.

\section{Work Plan}

The completed experimental artifacts and remaining presentation tasks are organized as follows.

\begin{table}[h]
\centering
\small
\begin{tabularx}{\textwidth}{p{0.24\textwidth}>{\raggedright\arraybackslash}X}
\toprule
\textbf{Stage} & \textbf{Activities and deliverables} \\
\midrule
Protocol and system & Fixed memory schema, LangGraph workflow, Qwen planner, ToolSandbox adapter, authoritative verifier, and six policies \\
Development and evaluation & 40 development and 120 held-out episodes; independent budget-one and budget-three runs with complete traces \\
Analysis & Offline forced-verification labels, clustered bootstrap, subgroup and cost analyses, and documented failure cases \\
Remaining deliverables & Review the report and contribution record, check key reproductions, and prepare slides, demo, recording, and timed rehearsal \\
\bottomrule
\end{tabularx}
\end{table}

\appendix
\section{Literature Search and Review Protocol}
\label{app:literature}

The related-work review protocol combines targeted database search with backward and forward citation chasing. Sources include ACL Anthology, arXiv, OpenReview, Google Scholar, and Semantic Scholar. Citations should be checked against the paper text and official metadata rather than search snippets alone.

\subsection{Search Queries}

\begin{itemize}
  \item \texttt{"agent memory" AND (counterfactual OR intervention OR ablation)}
  \item \texttt{"persistent memory" AND (verification OR provenance OR grounding)}
  \item \texttt{"memory selection" AND (downstream OR decision OR action)}
  \item \texttt{"LLM agent" AND ("stale memory" OR "memory interference"}\allowbreak\ \texttt{OR "memory poisoning")}
  \item \texttt{"value of information" AND ("LLM agent" OR "tool use")}
  \item \texttt{(planner OR replanning) AND memory AND counterfactual}
  \item \texttt{"agent memory" AND (benchmark OR evaluation) AND}\newline
  \hspace*{1em}\texttt{(multi-session OR long-horizon)}
\end{itemize}

Backward citation chasing starts from AMV-Guard, CMI, MemAudit, ConMem, CICL, TierMem, and MemoryArena; forward citation chasing is used where indexing supports it. Search records can be updated as recent preprints change.

\subsection{Inclusion and Exclusion Criteria}

Included works address at least one of the following: persistent memory in LLM agents; memory effects on downstream action; counterfactual memory intervention; provenance or verification; verification-cost allocation; or executable multi-session evaluation. Purely parametric knowledge editing, non-agentic long-context retrieval, and conversational recall without a decision component serve only as background. Peer-reviewed papers and preprints are distinguished explicitly.

\subsection{Evidence Extraction}

The evidence-extraction fields are: full citation and publication status; memory lifecycle stage; intervention target; whether the planner is re-executed; whether verification occurs before action; environment and task split; baselines; quality and cost metrics; claimed limitations; and available code or data. These fields support comparison with the closest available methods.

\bibliographystyle{plainnat}
\bibliography{references}

\end{document}
